\section{Introduction}
\label{sec:intro}

When a deployed language model is corrected, handed a fresh fact, or shown that its first
answer was wrong, we want it to listen. Yet a growing intuition in the community is that the
very training that makes models useful assistants---reinforcement learning (RL) from human
feedback and, more recently, RL on verifiable rewards---also makes them stubborn: the model
appears to decide where its turn is going before it starts writing, and then refuses to be
moved. If true, this would be a direct reliability and safety problem. The models we put in
front of users in high-stakes, multi-turn settings would be exactly the ones \emph{least}
able to absorb a correction once the conversation is underway---the ``lost-in-conversation''
failure.

\para{The hypothesis.} This paper tests a specific, three-part version of that intuition,
posed as a research question: do RL/post-trained LLMs \emph{pre-commit} to the answer they
will give---encoding it early and resisting revision when new information arrives---more than
their base counterparts? And can we \emph{rotate} which sub-fact of a decomposable question
is held certain to \emph{redistribute} the model's uncertainty onto the remaining parts? We
separate this into three claims. \textbf{(C1)} Pre-planning exists: the final answer is
decodable from internal state before it is emitted. \textbf{(C2)} RL increases
pre-commitment and reduces reactivity, both internally (the answer is encoded earlier) and
behaviorally (the model revises less). \textbf{(C3)} Rotating certainty redistributes
uncertainty---and RL models redistribute less.

\para{The gap.} Each ingredient has been studied, but never together on the comparison that
the hypothesis is actually about. Interpretability work firmly establishes C1: probes on the
prompt hidden state predict a response's length, answer, and correctness before the first
token~\citep{dong2025emergent,pal2023future,bigelow2025planning}. But that work does not
isolate RL---it never contrasts a matched base model against its RL-trained descendant.
Conversely, the behavioral literature on revision and sycophancy (\eg ``Try Again''
prompting, self-correction, FlipFlop challenges) measures reactivity directly but never ties
it to an internal pre-commitment probe, and it disagrees on the \emph{direction} of the
deficit: some report RL-induced \emph{under}-reaction and answer
repetition~\citep{ufo2025,kumar2024score}, others report widespread \emph{over}-reaction that
is not monotonic in RLHF~\citep{laban2023flipflop}. Finally, C3 has no precedent: no prior
work rotates what is held certain across a question to test whether uncertainty actually
spreads.

\para{Our approach.} We supply the missing contrast. We assemble a matched post-training
gradient of a \emph{single} model family---\qwenbase (base), \qweninst (SFT + RLHF/DPO), and
\rdistill (distilled from R1's RL reasoning traces)---so that architecture and tokenizer are
held fixed and post-training intensity is the variable. On the \emph{same} items we run three
experiments: (1) a two-turn \textbf{reactivity} test that injects content-free pushback or a
genuine hint and measures whether the model revises (C2b); (2) a \textbf{commitment-earliness
probe} that decodes the model's own final answer from its prompt hidden state and along the
generation timeline (C1, C2a); and (3) a \textbf{rotating-certainty} protocol on decomposable
\strat questions that measures how much answer entropy moves when one sub-fact is held certain
(C3).

\para{What we find.} The results invert the hypothesis. It is the \emph{base} model, not the
RL model, that most knows the answer it wants to give. The probe decodes the base model's
final answer from its prompt hidden state at \textbf{87\%} balanced accuracy (chance 25\%),
falling to 71\% for instruct and to near-chance \textbf{33\%} for the RL-reasoning model.
Behaviorally the base model is the most rigid: it revises only \textbf{2\%} of the time when
handed a genuinely informative hint, versus $\sim$47\% for both post-trained models
($p\approx10^{-27}$). Holding a sub-fact certain redistributes the base model's uncertainty
by $-0.32$ bits ($p\approx10^{-13}$) but leaves the reasoning model's entropy essentially
untouched ($p=0.84$). Post-training does not increase pre-commitment---it moves the
commitment point \emph{later} and makes the model reactive, at the cost (for plain RLHF) of
introducing the opposite failure, sycophantic over-reaction.

\para{Contributions.}
\begin{itemize}[leftmargin=*,itemsep=0pt,topsep=2pt]
    \item We provide the first matched base-vs-RL \emph{pre-planning-probe} contrast,
    decoding each model's own answer from its prompt representation, and show that internal
    pre-commitment \emph{decreases} monotonically with post-training (0.87 $\rightarrow$ 0.71
    $\rightarrow$ 0.33), with the commitment point moving later along generation.
    \item We tie this internal measure to behavior on two tasks, showing that the base model
    is the consistently rigid one while the two post-trained models exhibit the two opposite
    reactivity failures the literature flagged---instruct over-reacts (sycophancy), the
    reasoning model is selective on \arc but over-reacts on \gsm.
    \item We introduce and run the first test of \emph{rotating certainty}, showing that
    holding a sub-fact certain redistributes answer uncertainty---strongly for the base
    model, marginally for instruct, and not at all for the reasoning model.
    \item We release code, raw outputs, and figures for all three experiments, with logged
    parse-failure rates and bootstrap/Wilcoxon statistics, so the contrast is reproducible.
\end{itemize}

The rest of the paper reviews related work (\secref{sec:related}), details our models,
datasets, and metrics (\secref{sec:method}), presents the three experiments
(\secref{sec:results}), and discusses what the inverted result means and where it stops
(\secref{sec:discussion}).
